\chapter{Advanced Streams and File I/O}
\label{ch:streams-file-io}

\textit{Talking to the outside world --- files, formatted output, and the iostream hierarchy.}

C++ streams share a uniform model for console, file, and in-memory string I/O. This chapter covers: C-strings, \texttt{std::string}, the iostream class hierarchy, file I/O with \texttt{<fstream>}, stream states and error recovery, format manipulators from \texttt{<iomanip>}, binary I/O, string streams, stream internals, and custom manipulators.

%% ============================================================
\section{C-Strings: The Null-Terminated Array}
\label{sec:c-strings}

\begin{lstlisting}[language=C++, caption={C-string basics and null terminator}]
#include <cstring>
#include <iostream>

int main() {
    char name[6] = {'H', 'e', 'l', 'l', 'o', '\0'};
    const char* greeting = "World"; // Compiler adds '\0'

    std::cout << name << " " << greeting << "\n";
    std::cout << "Length: " << strlen(name) << "\n"; // 5
    return 0;
}
\end{lstlisting}

\begin{lstlisting}[language=C++, caption={Safer C-string manipulation}]
#include <cstring>
#include <cstdio>

int main() {
    char src[] = "Hello";
    char dst[20];

    strncpy(dst, src, sizeof(dst) - 1); // Bounded copy
    dst[sizeof(dst) - 1] = '\0';

    printf("Length: %zu\n", strlen(dst));
    printf("Compare: %d\n", strcmp("abc", "abd")); // negative
    return 0;
}
\end{lstlisting}

\textbf{Prefer \texttt{std::string}} in C++ code; use \texttt{strncpy}/\texttt{snprintf} only when interoperating with C APIs.

%% ============================================================
\section{\texttt{std::string}: The Safe Alternative}
\label{sec:std-string-ch10}

\begin{lstlisting}[language=C++, caption={std::string replaces C-strings}]
#include <string>
#include <iostream>
using namespace std;

int main() {
    string first = "John";
    string last  = "Wick";
    string full  = first + " " + last;

    cout << "Name: "   << full << "\n";
    cout << "Length: " << full.length() << "\n";

    const char* cstr = full.c_str(); // Null-terminated const pointer
    return 0;
}
\end{lstlisting}

%% ============================================================
\section{Basic Console I/O}
\label{sec:console-io}

\begin{center}
\begin{tabular}{|l|l|l|}
\hline
\textbf{Object} & \textbf{Direction} & \textbf{Description} \\
\hline
\texttt{std::cout} & Output & Standard output, buffered \\
\texttt{std::cin}  & Input  & Standard input (keyboard) \\
\texttt{std::cerr} & Output & Standard error, unbuffered \\
\texttt{std::clog} & Output & Standard error, buffered \\
\hline
\end{tabular}
\end{center}

\begin{lstlisting}[language=C++, caption={Console I/O and getline}]
#include <iostream>
#include <string>
using namespace std;

int main() {
    string line;
    cout << "Enter a sentence: ";
    getline(cin, line); // Reads until '\n', discards the newline
    cout << "You typed: " << line << "\n";
    return 0;
}
\end{lstlisting}

%% ============================================================
\section{Stream States and Error Recovery}
\label{sec:stream-states}

\begin{center}
\begin{tabular}{|l|l|l|}
\hline
\textbf{Flag} & \textbf{Access} & \textbf{Meaning} \\
\hline
goodbit & \texttt{good()} & No errors \\
eofbit  & \texttt{eof()}  & End of file reached \\
failbit & \texttt{fail()} & Non-fatal error (type mismatch) \\
badbit  & \texttt{bad()}  & Fatal error (hardware failure) \\
\hline
\end{tabular}
\end{center}

\begin{lstlisting}[language=C++, caption={Recovering from stream error state}]
#include <iostream>
using namespace std;

int main() {
    int age;
    cout << "Enter age: ";
    if (!(cin >> age)) {
        cout << "Invalid input.\n";
        cin.clear();            // Reset error bits
        cin.ignore(10000, '\n'); // Discard bad input
    }
    return 0;
}
\end{lstlisting}

%% ============================================================
\section{File I/O with \texttt{<fstream>}}
\label{sec:fstream}

\begin{lstlisting}[language=C++, caption={Writing and reading text files}]
#include <fstream>
#include <string>
#include <iostream>
using namespace std;

int main() {
    // Write
    ofstream out("save_data.txt");
    if (out.is_open()) {
        out << "PlayerLevel: 99\n";
        out << "Gold: 5000\n";
        out.close();
    }

    // Read line-by-line
    ifstream in("save_data.txt");
    if (!in.is_open()) { cerr << "File not found.\n"; return 1; }

    string line;
    while (getline(in, line))
        cout << "Read: " << line << "\n";
    in.close();
    return 0;
}
\end{lstlisting}

%% ============================================================
\section{Open Modes}
\label{sec:open-modes}

\begin{center}
\begin{tabular}{|l|l|}
\hline
\textbf{Flag} & \textbf{Meaning} \\
\hline
\texttt{ios::in}     & Open for reading \\
\texttt{ios::out}    & Open for writing (creates or truncates) \\
\texttt{ios::app}    & All writes go to end-of-file \\
\texttt{ios::ate}    & Open, seek to end \\
\texttt{ios::binary} & Binary mode --- no CRLF translation \\
\texttt{ios::trunc}  & Truncate existing file to zero \\
\hline
\end{tabular}
\end{center}

\begin{lstlisting}[language=C++, caption={Combining open modes}]
#include <fstream>
using namespace std;

int main() {
    ofstream log("app.log", ios::out | ios::app); // Append mode
    log << "Session started\n";
    log.close();

    fstream db("data.bin", ios::in | ios::out | ios::binary);
    if (db) { /* read/write binary data */ db.close(); }
    return 0;
}
\end{lstlisting}

%% ============================================================
\section{Binary File I/O}
\label{sec:binary-io}

\begin{lstlisting}[language=C++, caption={Binary write and read}]
#include <fstream>
#include <iostream>
using namespace std;

int main() {
    // Write
    {
        ofstream out("data.bin", ios::binary);
        int numbers[] = {10, 20, 30, 40, 50};
        out.write(reinterpret_cast<const char*>(numbers), sizeof(numbers));
    }
    // Read
    {
        ifstream in("data.bin", ios::binary);
        int buffer[5];
        in.read(reinterpret_cast<char*>(buffer), sizeof(buffer));
        for (int i = 0; i < 5; ++i) cout << buffer[i] << " ";
        cout << "\n";
    }
    return 0;
}
\end{lstlisting}

\begin{lstlisting}[language=C++, caption={Struct serialisation to binary file}]
#include <fstream>
#include <iostream>
using namespace std;

struct PlayerRecord { int level; float health; char name[32]; };

int main() {
    PlayerRecord rec = {99, 100.0f, "Arthur"};
    { ofstream out("player.bin", ios::binary);
      out.write(reinterpret_cast<const char*>(&rec), sizeof(rec)); }

    PlayerRecord loaded;
    { ifstream in("player.bin", ios::binary);
      in.read(reinterpret_cast<char*>(&loaded), sizeof(loaded)); }

    cout << loaded.name << " level=" << loaded.level << "\n";
    return 0;
}
\end{lstlisting}

%% ============================================================
\section{Stream Positioning}
\label{sec:stream-positioning}

\begin{lstlisting}[language=C++, caption={seekg, tellg, seekg with direction}]
#include <fstream>
#include <iostream>
using namespace std;

int main() {
    { ofstream out("test.txt"); out << "0123456789"; }

    ifstream in("test.txt");
    cout << "Position: " << in.tellg() << "\n"; // 0

    in.seekg(5);        // Seek to byte 5
    char c;
    in.get(c); cout << "At pos 5: " << c << "\n"; // '5'

    in.seekg(-3, ios::end); // 3 bytes before end
    in.get(c); cout << "3 from end: " << c << "\n"; // '7'

    in.close();
    return 0;
}
\end{lstlisting}

%% ============================================================
\section{Format Manipulators (\texttt{<iomanip>})}
\label{sec:manipulators}

\begin{lstlisting}[language=C++, caption={Integer base and prefix}]
#include <iostream>
#include <iomanip>
using namespace std;

int main() {
    int n = 255;
    cout << dec << n << "\n";              // 255
    cout << hex << n << "\n";              // ff
    cout << uppercase << hex << n << "\n"; // FF
    cout << showbase  << hex << n << "\n"; // 0XFF
    cout << oct << n << "\n";             // 0377
    cout << dec;                           // Reset to decimal
    return 0;
}
\end{lstlisting}

\begin{lstlisting}[language=C++, caption={Floating-point precision and notation}]
#include <iostream>
#include <iomanip>
using namespace std;

int main() {
    double pi = 3.14159265358979;
    cout << pi                                << "\n"; // 3.14159
    cout << fixed      << setprecision(2) << pi << "\n"; // 3.14
    cout << scientific << setprecision(4) << pi << "\n"; // 3.1416e+00
    cout << defaultfloat; // Reset
    return 0;
}
\end{lstlisting}

\begin{lstlisting}[language=C++, caption={Width, fill, and alignment}]
#include <iostream>
#include <iomanip>
using namespace std;

int main() {
    cout << "no setw: "   << 51 << "\n";
    cout << "setw(7): "   << setw(7) << 51 << "\n";             // "     51"
    cout << "setfill(*): " << setw(7) << setfill('*') << 51 << "\n"; // "*****51"
    cout << left  << setw(10) << setfill(' ') << "left"  << "|\n"; // "left      |"
    cout << right << setw(10) << "right" << "|\n"; // "     right|"
    cout << boolalpha << true << " " << false << "\n"; // true false
    return 0;
}
\end{lstlisting}

%% ============================================================
\section{String Streams (\texttt{<sstream>})}
\label{sec:sstream}

\begin{lstlisting}[language=C++, caption={ostringstream for in-memory serialisation}]
#include <sstream>
#include <string>
#include <iostream>
using namespace std;

int main() {
    ostringstream ss;
    ss << "The answer to everything is " << 42;
    string result = ss.str(); // "The answer to everything is 42"
    cout << result << "\n";
    return 0;
}
\end{lstlisting}

\begin{lstlisting}[language=C++, caption={istringstream for parsing}]
#include <sstream>
#include <string>
#include <iostream>
using namespace std;

int main() {
    string data = "Alice 30 3.14";
    istringstream ss(data);
    string name; int age; double value;
    ss >> name >> age >> value;
    cout << name << " age=" << age << " val=" << value << "\n";
    return 0;
}
\end{lstlisting}

\begin{lstlisting}[language=C++, caption={C++98 number-to-string idiom}]
#include <sstream>
#include <string>
using namespace std;

template<typename T>
string to_string_98(const T& val) {
    ostringstream ss; ss << val; return ss.str();
}

template<typename T>
T from_string_98(const string& s) {
    T result; istringstream ss(s); ss >> result; return result;
}

int main() {
    string s = to_string_98(42);
    int    n = from_string_98<int>("99");
    return 0;
}
\end{lstlisting}

%% ============================================================
\section{Stream Iterators}
\label{sec:stream-iterators}

\begin{lstlisting}[language=C++, caption={ostream\_iterator with algorithms}]
#include <algorithm>
#include <vector>
#include <iterator>
#include <iostream>
using namespace std;

int square(int x) { return x * x; }

int main() {
    int arr[] = {1, 2, 3, 4, 8, 16};
    vector<int> v(arr, arr + 6);

    copy(v.begin(), v.end(), ostream_iterator<int>(cout, " "));
    cout << "\n"; // 1 2 3 4 8 16

    transform(v.begin(), v.end(),
              ostream_iterator<int>(cout, " "), square);
    cout << "\n"; // 1 4 9 16 64 256
    return 0;
}
\end{lstlisting}

%% ============================================================
\section{Professional Insights: Stream Architecture}
\label{sec:stream-architecture}

\subsection{The \texttt{ios\_base} Class Hierarchy}

\begin{lstlisting}[language={}, caption={iostream class hierarchy}]
ios_base
└── basic_ios<charT>
    ├── basic_istream  → istream, ifstream, istringstream
    ├── basic_ostream  → ostream, ofstream, ostringstream
    └── basic_iostream → iostream, fstream, stringstream
\end{lstlisting}

\subsection{Redirecting Streams via \texttt{rdbuf}}

\begin{lstlisting}[language=C++, caption={Redirect cout to a file}]
#include <fstream>
#include <iostream>
using namespace std;

int main() {
    ofstream log("log.txt");
    streambuf* old = cout.rdbuf();  // Save original buffer
    cout.rdbuf(log.rdbuf());         // Redirect to file
    cout << "This goes to log.txt\n";
    cout.rdbuf(old);                 // Restore
    cout << "Back to screen\n";
    return 0;
}
\end{lstlisting}

\subsection{Custom Manipulators}

\begin{lstlisting}[language=C++, caption={User-defined stream manipulator}]
#include <iostream>
using namespace std;

ostream& tab(ostream& os)       { return os << "\t"; }
ostream& separator(ostream& os) { return os << "---\n"; }

int main() {
    cout << "Col1" << tab << "Col2" << "\n";
    cout << separator;
    return 0;
}
\end{lstlisting}

\subsection{Fast I/O: Decoupling C and C++ Streams}

\begin{lstlisting}[language=C++, caption={sync\_with\_stdio and cin.tie}]
#include <iostream>
using namespace std;

int main() {
    ios::sync_with_stdio(false); // 2-10x faster; do NOT mix with printf after this
    cin.tie(NULL);               // No auto-flush before cin reads

    int n;
    cin >> n;
    cout << n << "\n";
    return 0;
}
\end{lstlisting}

After calling \texttt{sync\_with\_stdio(false)}, mixing \texttt{printf}/\texttt{scanf} with \texttt{cout}/\texttt{cin} in the same program produces undefined output ordering.
